\chapter{Flat base change for the pushforward of a quasi-coherent sheaf ($i=0$)}
\label{chap:Cohomology_FlatBaseChange}
% archon:covers AlgebraicJacobian/Cohomology/FlatBaseChange.lean AlgebraicJacobian/Cohomology/FlatBaseChangeGlobal.lean
% SOURCE: [Stacks Project], Cohomology of Schemes, Section "Cohomology and
% base change, I", and Tag 02KH (read from references/stacks-coherent.tex and
% references/stacks-coherent.md).

\section{Setup and motivation}

This chapter establishes the $i=0$ (i.e.\ \(H^0\) / direct-image) case of flat
base change: the formation of the pushforward \(f_*\mathcal{F}\) of a
quasi-coherent sheaf commutes with flat base change on the target. It is the
foundational, fully self-contained slice of the cohomology-and-base-change
package that ultimately underlies relative representability statements; the
higher derived cases \(R^i f_*\) for \(i \geq 1\) are treated separately and are
out of scope here.

Throughout, \(f : X \to S\) is a morphism of schemes that is quasi-compact and
quasi-separated, and \(\mathcal{F}\) is a quasi-coherent \(\mathcal{O}_X\)-module.
A base change of \(f\) along an arbitrary morphism \(g : S' \to S\) is recorded
by the cartesian square obtained from the fibre product
\(X' = X \times_S S'\):
\[
  \begin{array}{ccc}
    X' & \xrightarrow{\;g'\;} & X \\[2pt]
    {\scriptstyle f'}\big\downarrow & & \big\downarrow{\scriptstyle f} \\[2pt]
    S' & \xrightarrow{\;g\;} & S
  \end{array}
\]
Here \(f' : X' \to S'\) is the base change of \(f\), and \(g' : X' \to X\) is the
first projection. We write \(\mathcal{F}' = (g')^*\mathcal{F}\) for the pullback
of \(\mathcal{F}\) to \(X'\).

\section{The canonical base-change map}


\section{The affine case}

\subsection{Locality of isomorphisms for quasi-coherent sheaf morphisms}

Before treating the affine computation we record three project-local supplements
to Mathlib. Mathlib provides the per-open criterion
(\(\mathrm{IsIso}\,\varphi \iff \forall U,\ \mathrm{IsIso}(\varphi_U)\)) and a
stalkwise isomorphism criterion for \(\mathrm{TopCat.Sheaf}\)-valued morphisms, but
it does not package the stalk-local or basis-local criterion at the level of
morphisms of \(\mathcal{O}_X\)-modules. The following lemmas bridge that gap; they
are exactly the locality tools used in the first reduction of the affine
base-change lemma below, where one checks the base-change map after restricting to
the affine opens of the base. As bespoke infrastructure they carry no external
citation and stand on their proof sketches.


\subsection{The affine computation}

The whole result rests on the following elementary computation in the affine
setting, where the base-change map is visibly an isomorphism because tensoring is
associative. Before stating it we isolate the one ingredient not yet available as
a ready-made result on which the affine reduction turns: the description of the
pushforward of a tilde-module along a morphism of affine schemes.

We first record three bespoke \(\Gamma\)-fragment supports that compute the global
sections of an affine pushforward; they carry no external citation and stand on
their proof sketches.


We record three further bespoke supports that drive the basic-open verification of
the affine-pushforward isomorphism below. They carry no external citation and stand
on their proof sketches.


\begin{remark}
  This single isomorphism discharges both inputs that the affine reduction below
  requires of the pushforward of a tilde-module.
  \begin{enumerate}
    \item \emph{The global-sections comparison.} Applying the global-sections
      (module-of-sections) functor to the isomorphism identifies the
      \(R\)-module of sections of
      \(\bigl(\operatorname{Spec}\varphi\bigr)_*\widetilde{M}\) with
      \(\operatorname{restr}_\varphi M\); this is exactly the \(\Gamma\)-fragment
      comparison used to recognise the section-level base-change map.
    \item \emph{Quasi-coherence of the pushforward.} The tilde of any module is
      quasi-coherent, and quasi-coherence of \(\operatorname{Spec}(R)\)-modules is
      closed under isomorphism; hence
      \(\bigl(\operatorname{Spec}\varphi\bigr)_*\widetilde{M}\) is quasi-coherent.
      In particular no separate ``pushforward preserves quasi-coherent'' theorem is
      needed for the affine lane.
  \end{enumerate}
  Because both inputs the affine reduction asks of the pushforward are corollaries
  of this single object isomorphism — the section-level module identification and
  the quasi-coherence of the pushforward — this lemma is the sole bespoke
  ingredient the affine lane requires.
\end{remark}

\subsection{The pullback companion}

The affine reduction below requires, alongside the pushforward dictionary, the
\emph{pullback} companion: the affine description of the inverse image of a
tilde-module. It is the affine case of the classical fact that the pullback of a
quasi-coherent sheaf is computed on modules by base change of scalars.

The pullback companion is built by uniqueness of left adjoints from the
per-object \(\Gamma\)-fragment comparison of
Lemma~\ref{lem:gammaPushforwardIso}, packaged as a natural isomorphism of
right-adjoint functors (Lemma~\ref{lem:gammaPushforwardNatIso}).


% SOURCE QUOTE PROOF: "The first assertion follows from the identification in
% Lemma \ref{lemma-compare-constructions}
% and the result of Modules, Lemma \ref{modules-lemma-restrict-quasi-coherent}."

\subsection{The affine base-change lemma and its remaining obligations}

The affine base-change lemma below proceeds \emph{directly on global sections}: it
reduces, via the locality criterion
Lemma~\ref{lem:modules_isIso_iff_affineOpens} and the two affine dictionaries
Lemmas~\ref{lem:pushforward_spec_tilde_iso} and~\ref{lem:pullback_spec_tilde_iso},
to two obligations that are isolated here as named blocks. Both are extracted from
the proof of Stacks ``Affine base change'', and both are absent from the pinned
Mathlib. The first is the \emph{locality reduction}: restricting to the affine opens
of \(S'\) reduces the isomorphism claim to the affine--affine model
\(S = \operatorname{Spec}(R)\), \(S' = \operatorname{Spec}(R')\),
\(X = \operatorname{Spec}(A)\), \(\mathcal{F} = \widetilde{M}\). The second is the
\emph{section-level identification}: in that affine--affine model, the global-sections
incarnation \(\Gamma(\alpha)\) of the base-change map, transported through the two
proved tilde dictionaries, is the canonical cancellation isomorphism
\(\operatorname{cancelBaseChange}\) — an isomorphism with no flatness hypothesis. This
is the honest cut of the affine close: locality on \(S'\), then a pure computation on
the module of global sections. No separate abstract adjoint-mate identification is
performed at the level of sheaves; the mate is unwound \emph{on sections}, where it
becomes \(\operatorname{cancelBaseChange}\) directly.


\subsection{The section-level base-change identity}

The section-level identification
Lemma~\ref{lem:pushforward_base_change_mate_cancelBaseChange} below is the single
highest-leverage obligation of the affine close. In the affine--affine model the
base-change map \(\theta\) is, after its domain and codomain are read through the two
proved tilde dictionaries, a concrete \(R'\)-linear map
\(R' \otimes_R M \to (R' \otimes_R A) \otimes_A M\), and the content of the lemma is
that this map is the inverse of the cancellation isomorphism
\(\operatorname{cancelBaseChange}\). Because the base-change map is \emph{defined} as the
transpose of the inner pushforward comparison \(\theta_{\mathrm{in}}\) under the
\((g^* \dashv g_*)\)-adjunction (\(g = \operatorname{Spec}\psi\)), its section-level value
is read off by the \emph{conjugate-counit calculus}: unfolding the transpose through the
counit form of the hom-set adjunction
(\(\operatorname{Adjunction.homEquiv\_counit}\)) factors \(\theta\) as
\(g^*(\theta_{\mathrm{in}})\) followed by the geometric counit \(\varepsilon_g\), and
conjugating this composite by the two tilde dictionaries
\(\Theta_{\mathrm{src}}\) (Lemma~\ref{lem:base_change_mate_domain_read}) and
\(\Theta_{\mathrm{tgt}}\) (Lemma~\ref{lem:base_change_mate_codomain_read}) replaces the
geometric counit by the \emph{algebraic} extension-of-scalars counit along \(\psi\). This
is the route the live formalization takes: the cancellation identity
Lemma~\ref{lem:pushforward_base_change_mate_cancelBaseChange} routes through the generator
trace Lemma~\ref{lem:base_change_mate_generator_trace}, which is a one-line corollary of the
section identity Lemma~\ref{lem:base_change_mate_section_identity}, whose single residual
obligation is the counit-conjugate coherence Lemma~\ref{lem:base_change_mate_gstar_transpose}.
The per-generator/element evaluation is \emph{not} available as a shortcut: evaluating the
opaque geometric composite on \(1 \otimes m\) produces no element-level normal form, so the
abstract conjugate calculus is the only tractable route. Two Mathlib results are recorded
first as dependency anchors.


% --- Supporting lemmas for the codomain read ---
% These helpers establish that pullback along a scheme isomorphism is an equivalence
% of module categories; the adjunction unit of an equivalence is an isomorphism,
% used in the codomain-read computation.


% --- Section-level identity, three pieces: regroup equiv (pure tensor algebra, native
%     R'-linear via Algebra.IsPushout.cancelBaseChange) ← section identity (Γ(θ) is the
%     R'-base change of the algebraic unit) ← IsIso corollary (the pinned leaf). ---


% --- Section-level identity, decomposed into three linked steps ---
%     The LHS unwinding of the base-change map on global sections splits at:
%       (1) lem:base_change_mate_unit_value  — the affine unit IS the algebraic unit;
%       (2) lem:base_change_mate_fstar_reindex — the pushforward pseudofunctor reindex;
%       (3) lem:base_change_mate_gstar_transpose — the (g^* ⊣ g_*) transpose on sections.


% --- Mathlib dependency anchors for the leg-reindex engine (mate / conjugate calculus) ---
% These three are the categorical tools the leg-reindex engine
% lem:pullbackPushforward_unit_comp rests on; they are supplied verbatim by Mathlib.


\begin{lemma}[The pullback--pushforward conjugate sends the pullback coherence to the pushforward coherence]
  \label{lem:conjugateEquiv_pullbackComp_inv_mathlib}
  \lean{AlgebraicGeometry.Scheme.Modules.conjugateEquiv_pullbackComp_inv}
  \mathlibok
  \textit{Provided by Mathlib.}
  Let \(f : X \to Y\) and \(g : Y \to Z\) be composable morphisms of schemes, and write
  \((h^{*} \dashv h_{*})\) for the pullback--pushforward adjunction on \(\mathcal{O}\)-modules
  attached to a morphism \(h\). Under the conjugation (mate) bijection between the composite
  adjunction \((g^{*}\dashv g_{*}) . (f^{*}\dashv f_{*})\) and the adjunction
  \(\bigl((f\circ g)^{*} \dashv (f\circ g)_{*}\bigr)\), the inverse pullback-composition coherence
  \((\mathrm{pullbackComp}_{f,g})^{\mathrm{inv}}\) is carried to the pushforward-composition
  coherence \((\mathrm{pushforwardComp}_{f,g})^{\mathrm{hom}}\):
  \[
    \operatorname{conj}\!\Bigl((g^{*}\dashv g_{*}) . (f^{*}\dashv f_{*}),\;
      \bigl((f\circ g)^{*}\dashv (f\circ g)_{*}\bigr)\Bigr)
      \bigl((\mathrm{pullbackComp}_{f,g})^{\mathrm{inv}}\bigr)
    \;=\; (\mathrm{pushforwardComp}_{f,g})^{\mathrm{hom}}.
  \]
\end{lemma}

% --- Mathlib dependency anchors for the conjugate-side re-encoding of the leg-reindex (Open Q2) ---
% These five categorical results are the tools of the conjugate-side (mate / Beck--Chevalley)
% proof route that replaces the abandoned direct-on-sections telescoping; all are supplied verbatim
% by Mathlib (CategoryTheory/Adjunction/Mates.lean and .../CompositionIso.lean).


% --- Seam 2, step (ii): Gamma-collapse of the transparent pushforward coherences ---
% These three project-bespoke atoms record that, on global sections over Spec R, the two
% pushforward-composition coherences and the pushforward-congruence coherence appearing in
% the inner composite are transparent (their section value is the identity / a presheaf
% transport), hence collapse under the global-sections functor.


% --- Seam 2, step (iii): the mate-unwinding crux, cut into five atomic links. ---
% These five lemmas decompose the step-(iii) obligation of
% lem:base_change_mate_fstar_reindex_legs (the mate-unwinding crux) into one
% mathematical move each. They are adjoint-mate calculus over the proved
% change-of-rings dictionaries; there is no external source. Composites
% are read left to right, as elsewhere in this chapter.


% --- Seam 2, step (iii): the mate-unwinding crux as a chain of five clean-term links. ---
% NOTE: These five lemmas re-express the step-(iii) telescoping of
% lem:base_change_mate_fstar_reindex_legs as a chain of equalities, each between two
% freshly stated terms over the same functor instance. This eliminates a nested-image vs
% composed-functor object diamond (\(G(H(M))\) vs \((H\circ G)(M)\)) that arises when
% distributing \(\Gamma\) over the composite. Each link is stated on its own clean composite
% where only one instance is in scope; the target chains the links and crosses the diamond once
% at the closing step. Throughout, \(e\) denotes the canonical square isomorphism,
% \(\iota_A = \mathrm{includeLeft}\), \(\Theta_{\mathrm{tgt}}\) the leg-parametrised codomain
% read (base_change_mate_codomain_read_legs), and \(\rho = \mathrm{base\_change\_mate\_inner\_value}\).


% --- Seam 2, step (iii) FINE re-break: the conjugate-side discharge of the
%     leg-reindex crux, atomised into one mathematical claim per lemma.
%     This chain cuts at the real conjugate-calculus seams. All blocks are adjoint-mate
%     calculus over the proved change-of-rings dictionaries and the conj-0 foundation
%     (lem:pullbackComp_eq_leftAdjointCompIso); there is no external source. Composites are
%     read left to right.


% NOTE: the `(g^* ⊣ g_*)`-transpose crux lem:base_change_mate_gstar_transpose cascades off
% this `_legs` identity by the SAME conjugate-side mechanism (its `Θ_src`/regroup/`Θ_tgt` move past the
% geometric factors is again a `conjugateEquiv` component identity, closed by the reassoc conjugate simp
% set + conjugateEquiv_pullbackComp_inv). It is a SEPARATE target.


% --- Seam A: the inner-value identity Γ_R(θ_in) = ρ. ---
% The (g^* ⊣ g_*)-transpose crux lem:base_change_mate_gstar_transpose needs the inner
% pushforward comparison θ_in, read on Spec R sections, to equal ρ : m ↦ (1 ⊗ 1) ⊗ m
% (lem:base_change_mate_inner_value_eq, Seam A). This identity is realised THROUGH the
% leg-parametrised reindex lem:base_change_mate_fstar_reindex_legs, instantiated at the literal
% projection legs via lem:base_change_mate_fstar_reindex: leg substitution → unit expansion
% (unitExpand) → four-factor Γ-distribution (gammaDistribute) → eCancel telescoping against the
% codomain read (the three one-cancellation atoms below, assembled in
% lem:base_change_mate_inner_eCancel_assemble) → Seam 1 → ρ.
% The inline alternative — distributing the bare (g')-unit AT the literal projection legs,
% without first passing to the leg-parametrised form — is WALLED: the leg pr_1 is only
% propositionally (not definitionally) equal to e ∘ Spec ιA, and the codomain read is a closed
% construction whose type bakes in the leg, so the unit expansion meets a dependent-motive
% obstruction at the literal leg. The proof must therefore pass to the leg-parametrised form
% first. The four-factor distribution (A-1) and the three eCancel atoms (A-2a/b/c) are correct
% and reused UNCHANGED; only the plumbing — via _legs, not inline — changes.
% The generator close (Seam B) follows as an independent lemma.


% --- The three one-cancellation atoms of the eCancel telescoping. ---
% Each is an independently-formalizable lemma whose proof is essentially a single move; together
% with the proved unit-distribution engine `lem:pullbackPushforward_unit_comp` they telescope the
% inner composite against the codomain read down to the lone affine `(Spec ιA)`-unit.


% --- Section-level identity (now a one-line chain of the three seams above). The base-change
%     map on sections is the R'-base change of the algebraic unit of extension of scalars;
%     reading it through the two tilde dictionaries pins it as the inverse of the regrouping
%     isomorphism. ---


\subsection{Tilde-transport direct proof of the section-level value}
\label{sec:fbc_tilde_transport_direct}

% This subsection gives a direct affine computation of the section-level value of
% the base-change map. It routes ENTIRELY through the single-step tilde dictionaries
% (Lemma~\ref{lem:pullback_spec_tilde_iso}, Lemma~\ref{lem:pushforward_spec_tilde_iso})
% and the tensor-product universal property; it does NOT depend on
% lem:base_change_mate_section_identity or lem:base_change_mate_gstar_transpose
% (the conjugate-calculus crux). Those conjugate scaffolding blocks remain in the
% chapter as a dictionary/record but are no longer on the critical path.

% NOTE (review iter-044): the "direct" bypass claimed above is ILLUSORY. iter-043 verified that
% the concrete affine square's inner unit is the `g'=pullback.fst` unit (no element normal form),
% so Gamma(alpha) must transit the same tilde dictionaries as the conjugate intertwining — this
% lemma collapses onto the SAME open keystone lem:base_change_mate_fstar_reindex_legs_conj. The
% Lean target `pushforward_base_change_mate_sections_direct` was NOT added (correctly). Treat this
% block as an abandoned route record, not a live critical path; do NOT dispatch a prover at it.

% SOURCE: [Stacks Project], Cohomology of Schemes, Lemma
% `lemma-affine-base-change' (\S\,Cohomology and base change, I), proof, the
% "boils down to the equality" step (read from references/stacks-coherent.tex,
% L927--938).
% SOURCE QUOTE PROOF: "We use Schemes, Lemma \ref{schemes-lemma-widetilde-pullback}
% to describe pullbacks and pushforwards of $\mathcal{F}$.
% Namely, $X' = \Spec(R' \otimes_R A)$ and
% $\mathcal{F}'$ is the quasi-coherent sheaf associated
% to $(R' \otimes_R A) \otimes_A M$.
% Thus we see that the lemma boils down to the
% equality
% $$
% (R' \otimes_R A) \otimes_A M = R' \otimes_R M
% $$
% as $R'$-modules."


% SOURCE QUOTE PROOF: "The vanishing of higher direct images is
% Lemma \ref{lemma-relative-affine-vanishing}.
% The statement is local on $S$ and $S'$. Hence we may
% assume $X = \Spec(A)$, $S = \Spec(R)$,
% $S' = \Spec(R')$ and $\mathcal{F} = \widetilde{M}$
% for some $A$-module $M$.
% We use Schemes, Lemma \ref{schemes-lemma-widetilde-pullback}
% to describe pullbacks and pushforwards of $\mathcal{F}$.
% Namely, $X' = \Spec(R' \otimes_R A)$ and
% $\mathcal{F}'$ is the quasi-coherent sheaf associated
% to $(R' \otimes_R A) \otimes_A M$.
% Thus we see that the lemma boils down to the
% equality
% $$
% (R' \otimes_R A) \otimes_A M = R' \otimes_R M
% $$
% as $R'$-modules."
% NOTE: This affine lemma is closed directly on global sections, via the two named
% obligations: lem:base_change_map_affine_local (locality reduction on the affine
% opens of S') and lem:pushforward_base_change_mate_cancelBaseChange (the section-
% level value of the base-change map is cancelBaseChange). Both proved affine tilde
% dictionaries -- lem:pushforward_spec_tilde_iso and lem:pullback_spec_tilde_iso --
% feed the second obligation. cancelBaseChange is in Mathlib with no flatness; the
% affine case carries no flatness hypothesis. No Čech infrastructure is used here.

\section{Flat base change for the pushforward}

The global statement is assembled from the affine lemma and the single homological
input that flatness preserves the finite equalizer computing \(H^0\). That input is
supplied verbatim by Mathlib, recorded here as a dependency anchor.


% NOTE: the former "A-linear equalizer presentation (FBC-B build-ahead)" subsection
% (3 blocks pinning gammaCoverRestrictScalars/gammaCoverResMapsALinear/gammaEqLocusIso) was RETIRED —
% replaced by the element-level route blocks (gammaModA/gammaResA/leftRes/rightRes/
% gammaTopEquivEqLocus/baseChangeGammaEquiv) in the FBC-B globalization subsection. The retired pins
% named never-built decls and created phantom frontier nodes.


% SOURCE QUOTE PROOF: "Having said this, part (1) follows from part (2). Namely,
% part (1) is local on $S'$ and hence we may assume $S$
% and $S'$ are affine. In other words, we have $S = \Spec(A)$
% and $S' = \Spec(B)$ as in (2).
% Then since $R^if_*\mathcal{F}$ is quasi-coherent
% (Lemma \ref{lemma-quasi-coherence-higher-direct-images}),
% it is the quasi-coherent $\mathcal{O}_S$-module associated to the
% $A$-module $H^0(S, R^if_*\mathcal{F}) = H^i(X, \mathcal{F})$
% (equality by
% Lemma \ref{lemma-quasi-coherence-higher-direct-images-application}).
% Similarly, $R^if'_*\mathcal{F}'$ is the quasi-coherent
% $\mathcal{O}_{S'}$-module associated to the $B$-module
% $H^i(X', \mathcal{F}')$. Since pullback by $g$ corresponds
% to $- \otimes_A B$ on modules
% (Schemes, Lemma \ref{schemes-lemma-widetilde-pullback})
% we see that it suffices to prove (2).
%
% Let $A \to B$ be a flat ring homomorphism.
% Let $X$ be a quasi-compact and quasi-separated scheme over $A$.
% Let $\mathcal{F}$ be a quasi-coherent $\mathcal{O}_X$-module.
% Set $X_B = X \times_{\Spec(A)} \Spec(B)$ and denote
% $\mathcal{F}_B$ the pullback of $\mathcal{F}$.
% We are trying to show that the map
% $$
% H^i(X, \mathcal{F}) \otimes_A B \longrightarrow H^i(X_B, \mathcal{F}_B)
% $$
% (given by the reference in the statement of the lemma)
% is an isomorphism.
%
% In case $X$ is separated, choose an affine open covering
% $\mathcal{U} : X = U_1 \cup \ldots \cup U_t$ and recall that
% $$
% \check{H}^p(\mathcal{U}, \mathcal{F}) = H^p(X, \mathcal{F}),
% $$
% see
% Lemma \ref{lemma-cech-cohomology-quasi-coherent}.
% If $\mathcal{U}_B : X_B = (U_1)_B \cup \ldots \cup (U_t)_B$ we obtain
% by base change, then it is still the case that each $(U_i)_B$ is affine
% and that $X_B$ is separated. Thus we obtain
% $$
% \check{H}^p(\mathcal{U}_B, \mathcal{F}_B) = H^p(X_B, \mathcal{F}_B).
% $$
% We have the following relation between the {\v C}ech complexes
% $$
% \check{\mathcal{C}}^\bullet(\mathcal{U}_B, \mathcal{F}_B) =
% \check{\mathcal{C}}^\bullet(\mathcal{U}, \mathcal{F}) \otimes_A B
% $$
% as follows from
% Lemma \ref{lemma-affine-base-change}.
% Since $A \to B$ is flat, the same thing remains true on taking cohomology.
%
% In case $X$ is quasi-separated, choose an affine open covering
% $\mathcal{U} : X = U_1 \cup \ldots \cup U_t$. We will use the
% {\v C}ech-to-cohomology spectral sequence
% Cohomology, Lemma \ref{cohomology-lemma-cech-spectral-sequence}.
% The reader who wishes to avoid this spectral sequence
% can use Mayer-Vietoris and induction on $t$ as in the proof of
% Lemma \ref{lemma-quasi-coherence-higher-direct-images}."
% NOTE: For the $i = 0$ case the argument is Čech-free: $H^0(X,\mathcal{F})$ is the
% sheaf-condition equalizer of a finite affine cover (Čech degree 0/1, NOT Čech
% cohomology -- no Leray, no spectral sequence), each term base-changes by the affine
% lemma lem:affine_base_change_pushforward, and flatness of $A \to B$ commutes
% $-\otimes_A B$ with that finite equalizer via lem:flat_preserves_equalizer_mathlib
% (Mathlib's LinearMap.tensorEqLocusEquiv). The quasi-separated case is the
% Mayer--Vietoris induction on the number of cover members named in the source quote.

\subsection{The \texorpdfstring{$A$}{A}-module \texorpdfstring{$\operatorname{eqLocus}$}{eqLocus} globalization (FBC-B)}

The separated-case argument is repackaged here at the level of \(A\)-modules, where
\(A = \Gamma(X,\mathcal{O}_X)\) is the ring of global sections of the structure sheaf, so that flat
base change becomes the Mathlib equivalence \(\operatorname{tensorEqLocusEquiv}\)
(Lemma~\ref{lem:flat_preserves_equalizer_mathlib}) applied to an honest
\(\operatorname{LinearMap.eqLocus}\). The blocks below construct that presentation concretely: every
module of sections \(\Gamma(U, M)\) is viewed as an \(A\)-module by restriction of scalars along the
structure-sheaf restriction \(A \to \Gamma(U,\mathcal{O}_X)\), the restriction maps of \(M\) become
\(A\)-linear, and the two product legs of the sheaf condition assemble into a pair of \(A\)-linear maps
whose equalizer locus is \(\Gamma(X, M)\). Throughout, \(M : X.\mathrm{Modules}\) is a sheaf of
\(\mathcal{O}_X\)-modules and \(U : \iota \to \mathrm{Opens}(X)\) is an open cover with
\(\bigvee_i U_i = \top\).

The bijectivity of the keystone presentation rests on two element-level sheaf axioms of the underlying
abelian-group sheaf \(M.\mathrm{presheaf}\), recorded as Mathlib dependency anchors.


\begin{lemma}[Sheaf gluing: a compatible family has a unique amalgamation]
  \label{lem:sheaf_existsUnique_gluing_mathlib}
  \lean{TopCat.Sheaf.existsUnique_gluing'}
  \mathlibok
  \textit{Provided by Mathlib.}
  Let \(\mathcal{G}\) be a sheaf on \(X\), \(U : \iota \to \mathrm{Opens}(X)\) a family covering an open
  \(V\), and \((s_i)_i\) a family of sections \(s_i \in \mathcal{G}(U_i)\) that is \emph{compatible}
  (the restrictions of \(s_i\) and \(s_j\) to \(U_i \cap U_j\) agree for all \(i, j\)). Then there is a
  unique \(s \in \mathcal{G}(V)\) whose restriction to each \(U_i\) is \(s_i\). This is the gluing half
  of the sheaf condition, stated element-wise.
\end{lemma}


\section{FBC-B: the global \texorpdfstring{$H^0$}{H0} flat-base-change isomorphism via the finite affine-cover equalizer (direct route)}

The blocks of the preceding subsection assemble into a \emph{direct} proof of the named global
isomorphism that bypasses the adjoint-mate keystone of the affine/mate route. The route is
\emph{{\v C}ech-cohomology-free}: \(H^0\) is a finite limit — the equalizer of a finite affine cover —
and the flat functor \(-\otimes_A B\) preserves finite limits, so the global comparison isomorphism
follows from per-chart affine base change \emph{without} the abstract conjugate/mate identification and
without the Mayer--Vietoris induction.

Concretely, the module-level core
\(B \otimes_A \Gamma(X,\mathcal{F}) \cong \operatorname{eqLocus}(\mathrm{id}_B\otimes\mathrm{leftRes},
\mathrm{id}_B\otimes\mathrm{rightRes})\) (Lemma~\ref{lem:fbcb_baseChangeGammaEquiv}) already commutes flat
base change past the finite \(H^0\) equalizer of a cover \(\{U_i\}\) of \(X\). The single remaining step
is to recognise the right-hand base-changed equalizer, chart by chart, as the equalizer-locus
presentation (Lemma~\ref{lem:fbcb_gammaTopEquivEqLocus}) of the global sections of the pulled-back module
over the base-changed cover \(\{(U_i)_B\}\) of \(X' = X_B\). The theorem below states that assembly.


